\documentclass[11pt]{article}

% * a(n) confirmed two ways up through a(22)                                          
% * a(n) computed using following techniques to cut search space up to a(35) or a(40),
%   depending on how strict we want to be                                             
% * split partial sums of a(n) into T(n,H)                                            
%    - exposes T(H, H) and T(k, H) for small k; this leads naturally to formulae for  
%      other regions of the triangle                                                  
% * computation of related sequences for symmetrical polyplets, those having holes    
\usepackage[margin=1in]{geometry}
\usepackage{amsmath,amssymb}
\usepackage{amsthm}
\usepackage{booktabs}
\usepackage{graphicx}
\usepackage{numprint}
\usepackage{microtype}
\usepackage[table]{xcolor}
\usepackage[hidelinks]{hyperref}
\usepackage{flafter}   % a float never appears before the text that introduces it
\usepackage{float}     % [H] pins tables/figures exactly where written
\usepackage[font=small,labelfont=bf,labelsep=period]{caption}
\setlength{\intextsep}{16pt plus 3pt minus 3pt}

\npthousandsep{,}
\npfourdigitsep
\newcommand{\num}[1]{\numprint{#1}}
\newcommand{\g}{\cellcolor{gray!15}}

\title{Polyplets}
\author{Jason Parker-Burlingham \texttt{<jasonp@panix.com>}}
\date{July 18, 2026}

\begin{document}
\maketitle

\begin{abstract} The king-lattice polyplet sequence \href{https://oeis.org/A006770}{A006770} has
been expanded through \[a(40) = \num{56749893611764175164545926946127}.\]

A new triangular sequence $T(n,H)$  is proposed that counts the fixed king polyplets of size $n$ and
height exactly $H\le{}n$. Clearly $a(n) = \sum_H\,T(n,H)$. Diagonals in the triangle ($H>n/2$) have
special forms such that a polynomial factor can be fixed from earlier rows and the other factor is
exponential in $n$ and $H$.

$a(1)$--$a(18)$ were already known; $a(19)$ was computed directly with Redelmeier enumeration;
$a(20)$--$a(23)$ were computed via transfer matrix; $a(24)$--$a(40)$ used earlier rows to fix
polynomial factors to compute the most expensive $T(n,H)$ by evaluation of the closed forms.

All $a(n)$ values for $n<40$ were double-checked; for $n\le{}22$, transfer matrix and Redelmeier
enumeration agree; for $n\le{}40$, a transfer-matrix method that uses coloring instead of directly
tracking connectivity confirms the $T(n,H\le{}19)$ cells.  All $a(n)$ values pass checks based on
Burnside congruences.

Every new $a(n)$ value has been checked by a second computation that shares no code; for $n\le{}22$
the confirmation comes from Redelmeier enumeration; for larger $n$ the confirmation is from a
transfer matrix that colors cells instead of tracking connectivity.

This work is completed with heavy assistance from Claude Code, especially the Fable model.
\end{abstract}

\section{Definitions}
\label{sec:definitions}

An \textbf{animal} is a finite, connected set of cells on a lattice; the \textbf{size} of an animal
is the number of cells that comprise it; the \textbf{height} of an animal is the height of its
bounding box.  In particular the \textbf{king polyplets} (or polyplets) that exist on the king
lattice where cells are connected by the 8 compass directions are considered here; the
rook-connected animals known as \textbf{polyominoes} are also relevant.

The polyplets can be counted with respect to the symmetries of the king lattice:
\begin{description}
\item[Fixed] polyplets are distinct up to translation (\href{https://oeis.org/A006770}{A006770});
\item[One-sided] polyplets are distinct up to translation and rotation
	(\href{https://oeis.org/A030233}{A030233});
\item[Free] polyplets are distinct up to translation, rotation and reflection
	(\href{https://oeis.org/A030222}{A030222});
\end{description}

Within the free polyplets, those that are \textbf{bilateral} have at least one reflection symmetry
(\href{https://oeis.org/A030234}{A030234}); those that are \textbf{asymmetric} have no reflective
symmetry (\href{https://oeis.org/A030235}{A030235}).

A polyplet may have one or more \textbf{holes}, polyominoes of unfilled cells enclosed by
king-connected occupied cells.  Thus a hole of size $1$ may be enclosed by a polyplet of size $4$,
and a domino of $2$ cells or two size-$1$ holes may be enclosed by a polyplet of size $6$.

$\mathbf{T(n,H)}$, $1\le{H}\le{n}$ is the number of fixed polyplets of size $n$ whose bounding box
is of height exactly $H$; the row sums produce $a(n)$, the number of fixed polyplets of size $n$
(\href{https://oeis.org/A006770}{A006770}).

\section{Results}
\label{sec:results}

% In narrative order, give:
% 1 a(n)
% 2 T(n,H), closed-form regions, diagonal law
% 3 One-sided, free, bilateral, asymmetric; A194596
% 4 Hole-classified counts
% 5 lambda bounds

\paragraph{Enumeration of the fixed polyplets} Table~\ref{tab:an} gives all known $a(n)$:

\begin{table}
\begin{tabular}{r r @{\quad}|@{\quad} r r}
\toprule
$n$ & $a(n)$ & $n$ & $a(n)$\\
\midrule
1           & \num{1}                & \textbf{21} & \num{6954084405510437}\\
2           & \num{4}                & \textbf{22} & \num{47255332844367680}\\
3           & \num{20}               & \textbf{23} & \num{321749260511448732}\\
4           & \num{110}              & \textbf{24} & \num{2194666793369310473}\\
5           & \num{638}              & \textbf{25} & \num{14994811325186658577}\\
6           & \num{3832}             & \textbf{26} & \num{102607513847014153892}\\
7           & \num{23592}            & \textbf{27} & \num{703126792093436034256}\\
8           & \num{147941}           & \textbf{28} & \num{4824589228356722264087}\\
9           & \num{940982}           & \textbf{29} & \num{33145129805782422061325}\\
10          & \num{6053180}          & \textbf{30} & \num{227969227118066423789154}\\
11          & \num{39299408}         & \textbf{31} & \num{1569631619081324581014300}\\
12          & \num{257105146}        & \textbf{32} & \num{10818203977457804418974036}\\
13          & \num{1692931066}       & \textbf{33} & \num{74631481980411777590683952}\\
14          & \num{11208974860}      & \textbf{34} & \num{515316838423862758858377704}\\
15          & \num{74570549714}      & \textbf{35} & \num{3561147281381782175236253062}\\
16          & \num{498174818986}     & \textbf{36} & \num{24629107617723857143962968288}\\
17          & \num{3340366308393}    & \textbf{37} & \num{170463735577007360431441250424}\\
18          & \num{22471158811164}   & \textbf{38} & \num{1180654489101178485738417779914}\\
\textbf{19} & \num{151609203011580}  & \textbf{39} & \num{8182864667276277865830132493466}\\
\textbf{20} & \num{1025573519362016} & \textbf{40} & \num{56749893611764175164545926946127}\\
\bottomrule
\end{tabular}
\caption{Fixed polyplets $a(n)$, $n=1,\dots,40$.  Terms $1$--$18$ match A006770; new terms have bold $n$.}
\label{tab:an}
\end{table}

\paragraph{Fixed polyplets by height} Table~\ref{tab:tnh} shows $T(n,H)$ for small $n$.  Regions of
this table have formulas: $T(n,n)$ is the number of polyplets with one cell per row; this
requirement forces $3^{n-1}$.  When $H>n/2$, cell placement choices are sufficiently constrained to
require $T(n,H) = P_{n-H}(n)\,3^{3H-2n-1}$ where the integer-valued $P_k$ has degree $k=n-H$ and the
leading coefficient is exactly $25^k/k!$.  $P_k$ is explicitly known for $k\le19$ after being fitted
against computed values of $T(n,H)$ and verified against values from later rows.  $P_k$ can be fixed
after computing the $3k$-th row of $T(n,H)$.

\begin{table}
\centering\footnotesize 
\setlength{\tabcolsep}{2.5pt}
\begin{tabular}{r | r r r r r r r r r r r r}
\toprule
$n$ & $H=1$ & 2 & 3 & 4 & 5 & 6 & 7 & 8 & 9 & 10 & 11 & 12\\
\midrule
 1 & \g1\\
 2 & 1 & \g3\\
 3 & 1 & \g10 & \g9\\
 4 & 1 & 27 & \g55 & \g27\\
 5 & 1 & 68 & \g248 & \g240 & \g81\\
 6 & 1 & 167 & 996 & \g1480 & \g945 & \g243\\
 7 & 1 & 406 & 3775 & \g7898 & \g7273 & \g3510 & \g729\\
 8 & 1 & 983 & 13837 & 39119 & \g47066 & \g32193 & \g12555 & \g2187\\
 9 & 1 & 2376 & 49676 & 185198 & \g278240 & \g241864 & \g133326 & \g43740 & \g6561\\
10 & 1 & 5739 & 175964 & 850650 & 1558399 & \g1631340 & \g1134865 & \g527094 & \g149445 & \g19683\\
11 & 1 & 13858 & 617789 & 3824200 & 8422286 & \g10311170 & \g8533676 & \g5001114 & \g2013255 & \g503010 & \g59049\\
12 & 1 & 33459 & 2156027 & 16922729 & 44372452 & 62471842 & \g59434367 & \g41336884 & \g21039624 & \g7487559 & \g1673055 & \g177147\\
\bottomrule
\end{tabular}
\caption{Fixed polyplets by height $T(n, H)$.  Shaded cells can be computed using formulas.}
\label{tab:tnh}
\end{table}

\paragraph{Non-fixed polyplets} Table~\ref{tab:onefree} and Table~\ref{tab:bisym} give new counts
for one-sided, free, bilateral and asymmetric polyplets of size $n$ as well as free polyplets that
are not polyominoes (\href{https://oeis.org/A194596}{A194596}).  Note that few polyplets are
polyominoes thanks to the higher polyplet growth rate, and that the one-sided count approaches
$a(n)/4$ and the free count approaches $a(n)/8$ as $n$ grows.

\begin{table}
\begin{tabular}{r | r r}
\toprule
$n$         &                         one-sided &                            free\\
\midrule
\textbf{18} &               \num{5617792259411} &             \num{2808898025438}\\
\textbf{19} &              \num{37902303297525} &            \num{18951156321090}\\
\textbf{20} &             \num{256393397108358} &           \num{128196711128365}\\
\textbf{21} &            \num{1738521118535082} &           \num{869260590388450}\\
\textbf{22} &           \num{11813833328245848} &          \num{5906916748001325}\\
\textbf{23} &           \num{80437315244081648} &         \num{40218657830371643}\\
\textbf{24} &          \num{548666699140240976} &        \num{274333350132318510}\\
\textbf{25} &         \num{3748702832087014098} &       \num{1874351417444073722}\\
\textbf{26} &        \num{25651878467205504568} &      \num{12825939237386225482}\\
\textbf{27} &       \num{175781698028751634291} &      \num{87890849023826053192}\\
\textbf{28} &      \num{1206147307126536440324} &     \num{603073653588819503780}\\
\textbf{29} &      \num{8286282451482505589564} &    \num{4143141225805220756331}\\
\textbf{30} &     \num{56992306779773178547095} &   \num{28496153390059675197641}\\
\textbf{31} &    \num{392407904770584270738465} &  \num{196203952385726340537921}\\
\textbf{32} &   \num{2704550994366217058008641} & \num{1352275497184284183618433}\\
\textbf{33} &  \num{18657870495104684580672401}\\
\textbf{34} & \num{128829209605977867230343869}\\
\bottomrule
\end{tabular}
\caption{One-sided and free polyplet counts.}
\label{tab:onefree}
\end{table}

\begin{table}
\begin{tabular}{r | r r r}
\toprule
$n$         &           bilateral &                      asymmetric &                 non-polyominoes\\
\midrule
\textbf{18} &       \num{3791465} &             \num{2808894233973} &             \num{2808705403386}\\
\textbf{19} &       \num{9344655} &            \num{18951146976435} &            \num{18950413696858}\\
\textbf{20} &      \num{25148372} &           \num{128196685979993} &           \num{128193840456415}\\
\textbf{21} &      \num{62241818} &           \num{869260528146632} &           \num{869249467327772}\\
\textbf{22} &     \num{167756802} &          \num{5906916580244523} &          \num{5906873556143637}\\
\textbf{23} &     \num{416661638} &         \num{40218657413710005} &         \num{40218489783363915}\\
\textbf{24} &    \num{1124396044} &        \num{274333349007922466} &        \num{274332695132618107}\\
\textbf{25} &    \num{2801133346} &       \num{1874351414642940376} &       \num{1874348860217028958}\\
\textbf{26} &    \num{7566946396} &      \num{12825939229819279086} &      \num{12825929238297403407}\\
\textbf{27} &   \num{18900472093} &      \num{87890849004925581099} &      \num{87890809870815114705}\\
\textbf{28} &   \num{51102567236} &     \num{603073653537716936544} &     \num{603073500077718909177}\\
\textbf{29} &  \num{127935923098} &    \num{4143141225677284833233} &    \num{4143140623183267694353}\\
\textbf{30} &  \num{346171848187} &   \num{28496153389713503349454} &   \num{28496151021712637626389}\\
\textbf{31} &  \num{868410337377} &  \num{196203952384857930200544} &  \num{196203943068019810549971}\\
\textbf{32} & \num{2351309228225} & \num{1352275497181932874390208} & \num{1352275460489267191905554}\\
\bottomrule
\end{tabular}
\caption{Counts of bilateral and asymmetric free polyplets, and the number of
non-polyomino free polyplets.  Note that
$\#\text{bilateral} + \#\text{asymmetric} = \#\text{free}$.}
\label{tab:bisym}
\end{table}

\paragraph{Holes} Polyplets with holes were tracked and counted during transfer
matrix enumeration by tracking the Euler characteristic through $n=18$.
Redelmeier enumeration checked these counts using flood-fill to verify counts
for $n\le14$.  Coounts are given in Table~\ref{tab:holes}

\begin{table}
\begin{tabular}{r | r r r r}
\toprule
$n$ & $k=0$ & 1 & 2 & 3\\
\midrule
 1 & 1\\
 2 & 4\\
 3 & 20\\
 4 & 109 & 1\\
 5 & 622 & 16\\
 6 & \num{3664} & 166 & 2\\
 7 & \num{22094} & \num{1456} & 42\\
 8 & \num{135609} & \num{11788} & 538 & 6\\
 9 & \num{843941} & \num{91300} & \num{5596} & 144\\
10 & \num{5310754} & \num{688034} & \num{52268} & \num{2082}\\
11 & \num{33724862} & \num{5091820} & \num{457754} & \num{24144}\\
12 & \num{215793158} & \num{37209939} & \num{3841776} & \num{248526}\\
13 & \num{1389673091} & \num{269452496} & \num{31290820} & \num{2374520}\\
14 & \num{8998648488} & \num{1937956658} & \num{249306876} & \num{21556520}\\
15 & \num{58548155506} & \num{13865448048} & \num{1953275736} & \num{188568632}\\
16 & \num{382526638033} & \num{98796744112} & \num{15103433472} & \num{1603972074}\\
17 & \num{2508473632910} & \num{701664375696} & \num{115556223626} & \num{13348948048}\\
18 & \num{16503616943998} & \num{4970078092354} & \num{876484089730} & \num{109173877126}\\
\bottomrule
\end{tabular}\\[2ex]
\begin{tabular}{r | r r r r r r r}
\toprule
$n$ & $k=4$ & 5 & 6 & 7 & 8 & 9 & 10\\
\midrule
 9 & 1\\
10 & 42\\
11 & 820 & 8\\
12 & \num{11482} & 263 & 2\\
13 & \num{135235} & \num{4808} & 96\\
14 & \num{1437090} & \num{67020} & \num{2186} & 22\\
15 & \num{14261808} & \num{804084} & \num{35086} & 808 & 6\\
16 & \num{134769521} & \num{8779911} & \num{465138} & \num{16410} & 314 & 1\\
17 & \num{1227552608} & \num{89826208} & \num{5490242} & \num{251104} & \num{7847} & 104\\
18 & \num{10866411352} & \num{876032532} & \num{59950610} & \num{3272674} & \num{137298} & \num{3460} & 30\\
\bottomrule
\end{tabular}
\caption{Count of fixed polyplets of size $n$ with exactly $k$ holes.}
\label{tab:holes}
\end{table}

% CLAUDE 2026-09-04, item 3. "checkable from Table 1": a reader who checks gets
%   a(40)^(1/40) = 6.2208 and a(40)/a(39) = 6.9352, not 7.11. 7.11 is an extrapolation
%   of the ratio sequence (paper/technical-report-gaps.md:33-38,
%   docs/main-paper-audit-2026-08-18.md:59-72).
%   Defensible from the table alone: a(m+n) >= a(m) a(n) (place one animal
%   diagonally off the corner of the other's box), so lambda = sup a(n)^(1/n) exists
%   and lambda >= a(40)^(1/40) = 6.2208  (results/concatenation-upper-bound.md:9).
%   Machine results, to cite rather than defend, if kept at all (your item 5):
%     ratio limit a(n+1)/a(n) -> lambda is a theorem, Madras 1999
%       (paper/technical-report-gaps.md:136-140)
%     lambda = 7.110(1), theta = -1.000(1) by differential approximants
%       (results/series-analysis-da.md:1)
%     6.543 <= lambda <= 9.3154, exact certificates, paper/L3-lambda-bounds.tex
\paragraph{Growth rate} $a(n)^{1/n} \to \lambda$ where $\lambda\approx7.11$, checkable from Table~\ref{tab:an}

\section{Methods}
\label{sec:methods}

\paragraph{Redelmeier enumeration} An implementation of Redelmeier enumeration was used to extend
$a(n)$ and to confirm the correctness of the counts for larger $n$.  As runtime and storage
increases with $a(n)$ only terms through $n=22$ could be confirmed with this software.

\paragraph{Transfer matrix method} When the space and time requirements for Redelmeier were reached,
polyplet classes were counted by height of the bounding box by freezing the left-hand-side of each animal
class in progress and considering how the polyplets could grow given constraints on the previous column,
whether the top and/or bottom of the bounding box had been reached, cell connectivity rules and the number
of cells still able to be placed.

\paragraph{Formulas} Formulas exist for $T(n,H)$ where $n=H$ and $H>n/2$.  $T(n,n)=3^{n-1}$
by simply choosing 1 of 3 positions to occupy in each row besides the first.

For $T(n, n-1)$, there are $n-1$ rows and $n-2$ junctions between rows where a choice is made for
the shape and placement of the next cell(s).  Single cells have exactly three options: left, right
and center of the previous cell.

To occupy $n$ cells in $n-1$ rows, one row must have two cells that both must be connected to the
rest of the polyplet:  a `joiner' that can be either a left-right domino or a `split' where two
cells are separated by exactly one empty cell (a larger gap could not be spanned by only one
occupied cell in the next row).  The joiner can be present at any of the $n-3$ interior rows.

For a domino joiner, one junction has 4 possibilities (rather than 3), and the other also has
4:  the domino may connect diagonally or orthogonally to the left or right of the previous cell,
and the subsequent cell may be similarly placed.  This means we have $16(n-3)\cdot{}3^{n-4}$ junction
possibilities.

%   ...X...  ...X...  ...X...  ...X...
%   .11....  ..22...  ...33..  ....44.
%   abcd...  .abcd..  ..abcd.  ...abcd

For a split joiner, when the joiner is exactly to the left and right of the previous cell, there are
five possibilities for the next row, otherwise there are four placements for the split that force a
single choice of the next row's cell.  That is 4 choices with one constrained subsequent choice, and
one choice with 5 choices, giving $(4\cdot{}1 + 1\cdot{}5)(n-3)\cdot{}3^{n-4}$ in total.

%   ...X...  ...X...  ...X...  ...X...  ...X...
%   1.1....  .2.2...  ..3.3..  ...4.4.  ....5.5
%   .a.....  ..a....  .abcde.  ....a..  .....a.

These combine to give us $25(n-3)\cdot{}3^{n-4}$ options for an interior joiner.

Next, account for the possibility that the polyplet starts or ends with a domino or a split; there
are five ways to start or end the polyplet, giving $2\cdot{}5\cdot{}3^{n-3} = 30\cdot{}3^{n-4}$
options to add to the interior joiner options:

%   ..11...  ..1.1..
%   .abcd..  ...e...

\[
    T(n, n-1) = 25(n-3)\cdot{}3^{n-4} + 30\cdot{}3^{n-4} = (25n - 45)\cdot{}3^{n-4}
\]

\section{Reproducibility}
\label{sec:reproduce}

\end{document}
